\chapter{Spike-train GLMs, resampling, and conditional hypotheses}\label{ch:glm}
\begin{readergoal}Build spike prediction models and valid null tests from conditional probabilities instead of selecting statistics first.\end{readergoal}
\section{A Poisson generalized linear model}
A GLM relates the expected count to covariates through a link function. For binned counts $n_k$ over exposure $\Delta_k$, one specification is
\[n_k\mid x_k\sim\mathrm{Poisson}(\mu_k),\qquad\log\mu_k=\log\Delta_k+\beta_0+x_k^\top\beta.\]
The offset $\log\Delta_k$ ensures that changing bin exposure has the correct interpretation in rate units. Under this model, $\lambda_k=\exp(\beta_0+x_k^\top\beta)$ is the conditional event rate. For sufficiently fine bins a Bernoulli model with appropriate link can also be used, but the likelihood assumptions differ.

\section{History filters and coupling}
A point-process GLM can augment predictors with spike-history bases or other units' recent spikes. Negative recent-self-history weights may express refractoriness; positive interactions may express predictive coupling. The estimated coefficients depend on basis choice, observation artifacts, unobserved common inputs, and regularization. Coupling filters are model parameters and not direct measurements of synaptic transmission.

\section{Model assessment rather than fit alone}
Inspect residual or time-rescaled-event structure, held-out log-likelihood, calibration, stability across bin widths, and sensitivity to covariates. Time-rescaling transforms integrated intensities between consecutive events; a correctly specified conditional-intensity model yields rescaled intervals with a known reference distribution. Failure suggests misspecification but does not identify the correct alternative automatically.

\section{Jitter and conditional resampling}
A resampling null model may fix the spike count inside each jitter window and randomize event locations within specified allowed regions. Exactness of conditional tests depends on the null distribution used for resampling. Generic ``jitter the spikes'' is not a statistical procedure until its preserved quantities, randomized quantities, boundary handling, and sampling distribution are fully specified. Conditional modeling clarifies the hypotheses supported by interval jitter and related tests \cite{Amarasingham2012}.

If $T_{\rm obs}$ is a test statistic and $T_1^*,\ldots,T_B^*$ are exchangeable null replicates under the test construction, a common Monte Carlo tail estimate is
\[\hat p=\frac{1+\sum_{b=1}^B\mathbf1\{T_b^*\ge T_{\rm obs}\}}{B+1}.\]
The plus-one convention avoids zero estimated $p$-values, but validity also requires the simulation to represent the intended null.

\begin{workedexample}[What is preserved?]
Suppose one window has three spikes. A conditional jitter method preserving this count cannot test whether the process increased its number of spikes in that window; it can test aspects of fine timing within the window. Rejection must be interpreted relative to preserved coarse structure, not as evidence for a particular synapse.
\end{workedexample}
\begin{exercises}
\item Derive the score of the Poisson log-likelihood with respect to $\beta$.
\item State the null hypothesis and exact conditioning in a 20-ms interval jitter test.
\item Explain why randomizing every timestamp independently across a session may destroy behavioral modulation.
\item Design a split-trial validation plan that avoids training and testing on overlapping windows from the same trial.
\end{exercises}